\section{From Models to Physical AI Systems}
\label{sec:systems}

The roadmap from LLM-based world knowledge to Physical AI ultimately requires deployed systems that can act in real or interactive environments. 
This section discusses how current models connect to policy learning, embodied agents, robotics, frontier systems, and evaluation.

\subsection{Policy Learning and Embodied Agents}
Policy learning connects perception, prediction, and action. 
In Physical AI, policies may be learned from demonstrations, reinforcement learning, offline datasets, synthetic trajectories, or world-model-based simulation. 
Embodied agents integrate high-level reasoning, multimodal perception, action execution, and feedback to operate in physical or simulated environments.

\subsection{Robotics and the Embodiment Gap}
Robotics provides a critical testbed for Physical AI, but it also exposes major gaps. 
Different robots have different bodies, sensors, action spaces, dynamics, and safety constraints. 
A policy that works for one embodiment may not transfer to another. 
This embodiment gap limits the generalization of current VLA and world-model-based approaches.

\subsection{Frontier Systems and Closed Models}
Recent frontier systems demonstrate rapid progress toward agentic and physical capabilities, including LLM agents, computer-use agents, robotics foundation models, and world-model platforms. 
However, many of these systems are closed-source or partially disclosed, making them difficult to evaluate reproducibly. 
A survey of Physical AI must therefore consider not only academic models, but also product-level systems that shape the field while lacking standard academic references.

\subsection{Benchmarks and Evaluation}
Evaluation remains a central challenge. 
Open-loop language or vision benchmarks are insufficient for Physical AI, which requires closed-loop interaction, physical consistency, task completion, robustness, safety, and generalization across environments and embodiments. 
Future benchmarks should evaluate not only what a model knows, but how reliably it can ground that knowledge into perception, prediction, planning, and action.